%% ===========================================================================
\FloatBarrier
\section{Particular cases}
\label{sec:limits}
%% ===========================================================================

The limiting cases are obtained by suppressing individual constitutive
contributions in the constitutive model. Each
switch acts on one term of the energy: \eqref{eq:Eel} for the elastic part,
\eqref{eq:Epiezo} for the piezoelectric and dielectric parts and
\eqref{eq:Eflexo} for the flexoelectric one. Each reduction is defined by the corresponding constitutive terms. These limits provide reference cases for the numerical verification in Sec.~\ref{sec:verification}.
All reductions follow from the stated enthalpy;
the classical limits are cited at the point of use.

\FloatBarrier
\subsection{Classical piezoelectricity}
\label{sec:limit_classical}

If the averaging window shrinks to a point, $M_{mn}\to0$, the gradient energy
\eqref{eq:Egrad} vanishes and the enthalpy reduces to
\begin{equation}
  \begin{aligned}
  \Pi \to {}&\tfrac12\int_\Omega \varepsilon_{ij} C_{ijkl}\varepsilon_{kl}
    \,\mathrm{d}A - \int_\Omega e_{kij}E_k\varepsilon_{ij}\,\mathrm{d}A\\
    &- \tfrac12\int_\Omega\kappa_{ij}E_iE_j\,\mathrm{d}A
    - \int_\Omega \mu_{kijl}E_k\eta_{ijl}\,\mathrm{d}A .
  \end{aligned}
  \label{eq:limitM}
\end{equation}
which is not yet classical piezoelectricity. By \eqref{eq:tau} the
higher-order stress retains the flexoelectric part $-\mu_{kijm}E_k$, so a model
with a window of zero size and a nonzero flexoelectric coefficient keeps a
higher-order stress, a double traction and a fourth-order character. The
classical limit is therefore the double limit $M_{mn}\to0$ \emph{and}
$\mu_{kijl}\to0$, in which the higher-order double-traction condition drops out of the formulation. What remains is the classical
piezoelectric strip: the weak form keeps only the elastic, piezoelectric and
dielectric terms, the displacement field needs $C^0$ continuity only, and the
coupling factor is the purely piezoelectric one. Thus, with flexoelectricity suppressed, the window size alone switches the
strain-gradient stiffness on or off without changing the remaining terms.

\FloatBarrier
\subsection{Purely piezoelectric response}
\label{sec:limit_piezo}

With the flexoelectric tensor set to zero the electric displacement
\eqref{eq:D} loses its strain-gradient contribution and the model reduces to
the linear piezoelectric strip of the plane-stress reduction \eqref{eq:Cp},
with the strain-gradient stiffness retained. For the length-poled extensional configuration, the coupling factor is $k^2=0.5727$ and is unchanged
by the switch, because that configuration collects the piezoelectric charge
only. On the bending configuration the piezoelectric charge of the mode
integrates to zero over a face electrode, and the piezoelectric-only coupling
factor of that mode falls to $1.0\times10^{-10}$, which is at the numerical round-off level of the secular calculation. The different coupling factors reflect the charge collected by the
respective electrode configurations.

\FloatBarrier
\subsection{Purely flexoelectric response}
\label{sec:limit_flexo}

Removing the piezoelectric coupling instead leaves the flexoelectric bending
response of the same strip. Over the thickness range
$\SI{10}{\micro\meter}\to\SI{1}{\micro\meter}$ the flexoelectric coupling of
that configuration then rises from $1.76\times10^{-4}$ to
$5.5\times10^{-3}$, a factor $31$. With the piezoelectric terms present the
same sweep gives a rise from $1.67\times10^{-4}$ to $4.0\times10^{-3}$, a
factor $24$. The piezoelectric contribution to collected charge vanishes for this
configuration, although the corresponding stiffness contribution remains: they change the elastic stiffness of the
mode that carries the gradient, and the difference between $24$ and $31$ is
the source of the difference. The factor $24$ refers to the model
with both couplings retained.

\FloatBarrier
\subsection{Classical Kelvin--Voigt and undamped limits}
\label{sec:limit_memory}

For $\alpha=1$ the local fractional material law becomes the classical
Kelvin--Voigt law, $\sigma=C(\varepsilon+\tau\dot\varepsilon)$,
with constitutive modulus $C$ in pascals. For a single proportional-memory
mode with modal stiffness $k_m$ and mass $m_m$, the dashpot coefficient
is $c_m=\tau k_m$, and the damping ratio is
\begin{equation}
  \zeta = \frac{c_m}{2\sqrt{k_m m_m}}
         =\frac{\tau k_m}{2\sqrt{k_m m_m}}
         =\frac{\omega_m\tau}{2},\qquad
  \omega_m=\sqrt{k_m/m_m},
  \label{eq:limitKV}
\end{equation}
where $k_m$ and $m_m$ are modal quantities, distinct from the material
modulus and the window second-moment tensor. This exact relation applies
when the memory stiffness is proportional to the same modal elastic
stiffness; it is not an exact per-mode formula for the non-proportional
mechanical-only memory operator.
The implementation reproduces it: at $\omega\tau=0.3$ the simulated
ring-down returns $\zeta_{\rm eff}=0.15003$ against the exact
$0.15000$, a difference of $0.02\,\%$. For $\tau\to0$ the memory term of
\eqref{eq:kv} vanishes and the time-discrete problem \eqref{eq:step}
becomes undamped. Its eigenvalues are those of the chosen
\emph{electrode-conditioned} elastic pencil: $K_{oc}$ for the reported
open-circuit ring-down and $K_{sc}$ for the corresponding short-circuit
problem. Between these two limits the model is genuinely fractional: the
damping ratio is no longer proportional to $\omega\tau$, and the choice of stiffness-memory matrix changes $\zeta_{\rm eff}$ by a factor $0.431$
at $\alpha=1$ for the specified device.

\FloatBarrier
\subsection{Recovery of the one-scale gradient theory}
\label{sec:limit_fem1}

Finally, the one-scale gradient theories are recovered as special choices of
$M$. For an ellipsoidal window $M_{mn}=L_{mn}/5$, so
$\tfrac14M_{mn}=\tfrac1{20}L_{mn}$ and the tensorial strain-gradient energy of
the one-scale theory~\cite{fem1} is reproduced exactly, with the cuboid,
ellipse and ellipsoid windows giving their respective scalar second-moment
coefficients; setting $M_{mn}=\ell^{2}\delta_{mn}$ reduces
the same expression to the isotropic Aifantis form~\cite{Aifantis1992}.
Dropping the memory and the flexoelectric tensor leaves a strain-gradient
piezoelectric strip; the classical strip additionally requires $M_{mn}=0$.
In the present formulation, fractional memory acts on the reduced modal
coordinates and the gradient stiffness depends on the window second-moment
tensor rather than on a scalar length. The same boundary and electrode conditions are imposed in these reductions.

\FloatBarrier
